\begin{figure}[htbp]
\centering
\includegraphics[width=\textwidth]{half_twisted_rectangular_tube_geometry.png}
\caption{The half-twisted rectangular tube and its two face-pair strips.
  Left: the complete surface, which is a closed orientable torus. Middle:
  the wide strip, $p=(s,b)$ with $s\in[-a,a]$. Right: the narrow strip,
  $p=(a,s)$ with $s\in[-b,b]$. Each strip covers its face pair once over
  $\theta\in[0,4\pi)$, and each is an annulus, not a M\"obius strip. The
  red and the blue curve are the two closed corner curves where the
  strips meet. The cross-section sides are straight, but the faces they
  sweep are not planar.}
\label{fig:tubegeom}
\end{figure}

\subsection{The bands and the routed rows}

\begin{figure}[htbp]
\centering
\includegraphics[width=0.92\textwidth]{half_twisted_rectangular_tube_narrow_band.png}
\caption{The two narrow bands at $h=1/22$, and the nodes the glue
  rotates. Blue dots are the outer band of branch~A, orange rings the
  outer band of branch~B, and magenta the rows that are routed across a
  corner curve. A node that lies in both bands carries both markers,
  because it carries two independent unknowns.}
\label{fig:tubeband}
\end{figure}

The routed rows are not a thin shell. Table~\ref{tab:tubesize} gives the
counts. At every level about one unknown in five is a row whose value is
supplied by the glue rather than by its own branch. That is the reason
the junction term in Equation~\eqref{eq:bound} is not a local
perturbation on this surface: there are two closed junction curves, and
every one of their neighbourhoods is routed.

\begin{table}[htbp]
\centering
\caption{Problem size and solver behaviour, half-twisted rectangular
  tube. Routed rows are summed over both branches and both corner
  curves. The row-sum error is $\max_i |\sum_j E_{ij} - 1|$.}
\label{tab:tubesize}
\small
\begin{tabular}{@{}rrrrrrr@{}}
\toprule
$h$ & unknowns A & unknowns B & routed rows & GMRES iter
    & rel.\ residual & row-sum error \\
\midrule
0.05000 &  32\,393 & 17\,391 & 21\,029 &  722 & 9.7e-11 & 1.3e-15 \\
0.04000 &  49\,040 & 25\,556 & 26\,686 &  840 & 9.9e-11 & 1.6e-15 \\
0.03333 &  68\,599 & 34\,963 & 31\,930 & 1203 & 9.9e-11 & 1.6e-15 \\
0.02778 &  96\,784 & 48\,138 & 37\,808 & 1544 & 9.9e-11 & 1.6e-15 \\
0.02273 & 141\,729 & 69\,344 & 46\,307 & 2276 & 1.0e-10 & 1.6e-15 \\
\bottomrule
\end{tabular}
\end{table}

\subsection{Convergence of the $\Linf$ error}

\begin{figure}[htbp]
\centering
\includegraphics[width=\textwidth]{half_twisted_rectangular_tube_rotation_convergence.png}
\caption{Left: the surface $\Linf$ error of both rotations, with $O(h)$
  and $O(h^2)$ guides. The two curves lie on top of one another. The
  error starts between the guides and bends towards $O(h)$ as $h$ falls.
  Right: the worst $d_{k2}$ rotation angle error, against an $O(h^2)$
  guide. It falls faster than second order.}
\label{fig:tubeconv}
\end{figure}

\begin{table}[htbp]
\centering
\caption{Surface $\Linf$ error of $u$, combined over both strips,
  half-twisted rectangular tube with $R=1$, $a=0.45$, $b=0.20$. Rates
  use the actual $h$ ratios.}
\label{tab:tubeconv}
\small
\begin{tabular}{@{}r rr rr r@{}}
\toprule
 & \multicolumn{2}{c}{exact rotation}
 & \multicolumn{2}{c}{$d_{k2}$ rotation}
 & $d_{k2}$ worst \\
\cmidrule(lr){2-3}\cmidrule(lr){4-5}
$h$ & error & rate & error & rate & angle error (rad) \\
\midrule
1/20 & 1.4041e-01 & --   & 1.4018e-01 & --   & 1.3967e+00 \\
1/25 & 8.5293e-02 & 2.23 & 8.5201e-02 & 2.23 & 6.7693e-01 \\
1/30 & 6.3814e-02 & 1.59 & 6.3876e-02 & 1.58 & 2.6388e-01 \\
1/36 & 5.4570e-02 & 0.86 & 5.4621e-02 & 0.86 & 1.3056e-01 \\
1/44 & 4.4203e-02 & 1.05 & 4.4186e-02 & 1.06 & 6.2282e-02 \\
\midrule
\multicolumn{1}{@{}l}{fitted}
     & \multicolumn{2}{c}{\textbf{1.429}}
     & \multicolumn{2}{c}{\textbf{1.426}}
     & 4.04 \\
\bottomrule
\end{tabular}
\end{table}

\paragraph{The fitted rate is 1.43, and it is not an asymptotic order.}
Table~\ref{tab:tubeconv} gives a least-squares slope of $1.429$ for the
exact rotation. That number alone is misleading, because the
level-to-level rates fall steadily:
\begin{equation*}
  2.23, \quad 1.59, \quad 0.86, \quad 1.05 .
\end{equation*}
The solve starts close to second order on the coarsest pair and settles
near first order on the finest pairs. Figure~\ref{fig:tubeconv} shows the
same thing: the error curve starts below the $O(h^2)$ guide and bends
towards the $O(h)$ guide. The fitted $1.429$ is an average over a range
in which the rate is still changing, not the order the scheme will hold
at smaller $h$. The honest statement is that the observed behaviour is
consistent with a first-order junction term taking over from a
second-order interior term, which is what Equation~\eqref{eq:bound}
predicts when $D$ and $J_2$ do not vanish.

Unlike the ellipsoid, this surface has no grid-alignment scatter: the
error is monotone in $h$ at every level. The two junction curves are not
plane circles, and the glue angle varies along them, so the signed
junction residual does not add up to a near cancellation that depends on
grid position.

\paragraph{Cross-check against the recorded run.} The notes in
\texttt{angle\_rotation\_glue.md} record a run at
$h\in\{1/25, 1/30, 1/40\}$. The three levels this study shares with it
reproduce it to every printed digit: $8.5293\mathrm{e}{-2}$ and
$6.3814\mathrm{e}{-2}$ for the exact rotation, $8.5201\mathrm{e}{-2}$ and
$6.3876\mathrm{e}{-2}$ for $d_{k2}$, and the level rate $1.59$ for the
pair $1/25 \to 1/30$.

\subsection{The estimated rotation costs nothing here}

The two rotations are indistinguishable in the solution. The largest
relative difference over all five levels is $0.16$ per cent, at the
coarsest level, and at the finest level the $d_{k2}$ solve is marginally
the \emph{better} of the two ($4.4186$ against
$4.4203 \times 10^{-2}$). That is not a real advantage; it is the size
of the difference.

This holds although the $d_{k2}$ rotation is a poor estimate on the
coarse grids. Its worst angle error is $1.40$~rad at $h=1/20$, which is
$80^\circ$ on the worst row. The error falls at a fitted rate of $4.04$,
with level rates $3.25$, $5.17$, $3.86$ and $3.69$, reaching
$6.2\times10^{-2}$~rad at the finest level.

\paragraph{How many rows need the fallback probe.} The $d_{k2}$ source
estimate is re-sampled from the deterministic outward probe
$c + h\,\eta_s$ whenever the reflected ray is unusable. The counts are
substantial and they grow with refinement:

\begin{center}
\small
\begin{tabular}{@{}rrrr@{}}
\toprule
$h$ & probed rows, branch A & probed rows, branch B & sign reversals \\
\midrule
1/20 & 2972 & 1789 & 7 \\
1/25 & 3100 & 2248 & 0 \\
1/30 & 3000 & 2609 & 0 \\
1/36 & 3360 & 3107 & 0 \\
1/44 & 3973 & 3854 & 0 \\
\bottomrule
\end{tabular}
\end{center}

So roughly one routed row in six falls back to the probe. The probe is
placed with the \emph{analytic} conormal, which is why this variant is
not a geometry-free estimator, as the script states. The seven sign
reversals at the coarsest level are rows where the long reflected sample
misjudged which side of the edge it was on.

\subsection{The traces agree at nearly second order}

The two branches agree with each other on the corner curves far better
than either agrees with the exact solution:

\begin{center}
\small
\begin{tabular}{@{}rrr@{}}
\toprule
$h$ & $\max|u_A-u_B|$, upper curve & $\max|u_A-u_B|$, lower curve \\
\midrule
1/20 & 2.879e-02 & 2.765e-02 \\
1/25 & 1.610e-02 & 2.123e-02 \\
1/30 & 1.219e-02 & 1.288e-02 \\
1/36 & 7.763e-03 & 1.144e-02 \\
1/44 & 6.412e-03 & 7.186e-03 \\
\midrule
\multicolumn{1}{@{}l}{fitted} & \textbf{1.93} & \textbf{1.71} \\
\bottomrule
\end{tabular}
\end{center}

The trace discrepancy converges at $1.93$ and $1.71$, against $1.43$ for
the surface error itself, and this is the same pattern the ellipsoid
shows. The glue transmits consistently across the junction at close to
second order. What limits the solve is not the transmission.
